Let's take a look about the things seen so far.

\hspace{3.5pt} \textbf{CP}.
\begin{itemize}
    \renewcommand{\labelitemi}{+}
    \item A generic and complete approach with a focus on constraints and feasibility. We can apply it to model any optimization problem and giving to it enough time it will explore the whole search tree, finding a feasible solution.
    \item Easy modelling and control of search.
\end{itemize} \vspace{3.5pt}
\begin{itemize}
    \renewcommand{\labelitemi}{-}
    \item Cannot scale to large optimization problems. The search tree can be quite huge, making the search process very challenging.
    \item Poor in optimization with loose bounds on the objective function.
\end{itemize} 

\hspace{3.5pt} \textbf{Heuristic search}.
\begin{itemize}
    \renewcommand{\labelitemi}{+}
    \item Scales to large optimization problems. For instance, local search or meta-heuristics are the kind of techniques that do no search over the whole search tree, but in portions of it.
    \item Effective in finding good-quality solutions quickly.
\end{itemize} \vspace{3.5pt}
\begin{itemize}
    \renewcommand{\labelitemi}{-}
    \item Finding an initial solution and exploring a large neighborhood can be challenging.
    \item Neighborhoods are problem-specific.
\end{itemize} 
Why don't we combine both of them, trying to explore exhaustively the search tree as CP and improving the solution in a more effective way like heuristics?